\documentclass{article}
\usepackage{amsmath, amssymb, amsthm}
\usepackage{hyperref}
\usepackage{geometry}
\geometry{margin=1in}

\title{Erd\H{o}s Problem \#860}
\date{Last updated: 2025-08-31}
\author{Source: erdosproblems.com}

\begin{document}
\maketitle

\noindent\textbf{Status:} OPEN \quad \textbf{No prize}

\noindent\textbf{Tags:} number theory, primes

\medskip

\noindent\textbf{Problem Statement:}

Let $h(n)$ be such that, for any $m\geq 1$, in the interval $(m,m+h(n))$ there exist distinct integers $a_i$ for $1\leq i\leq \pi(n)$ such that $p_i\mid a_i$, where $p_i$ denotes the $i$th prime. Estimate $h(n)$.




A problem of Erd\H{o}s and Pomerance \cite{ErPo80}, who proved that\[h(n) \ll \frac{n^{3/2}}{(\log n)^{1/2}}.\]Erd\H{o}s and Selfridge proved $h(n)&gt;(3-o(1))n$, and Ruzsa proved $h(n)/n\to \infty$. This is discussed in problem B32 of Guy's collection \cite{Gu04}. See also [375] .




References

[ErPo80] P. Erd\H{o}s and C. Pomerance, Matching the natural numbers up to $n$ with distinct multiples of another interval . Indigationes Math. (1980), 147-151.

[Gu04] Guy, Richard K., Unsolved problems in number theory . (2004), xviii+437.

\noindent\textbf{Related OEIS sequences:} A048670, A058989


\bigskip
\noindent\small{Source: \url{https://www.erdosproblems.com/860}}
\end{document}
